\let\negmedspace\undefined
\let\negthickspace\undefined
\documentclass[journal]{IEEEtran}
\usepackage[a5paper, margin=10mm, onecolumn]{geometry}
\usepackage{lmodern} % Ensure lmodern is loaded for pdflatex
\usepackage{tfrupee} % Include tfrupee package

\setlength{\headheight}{1cm} % Set the height of the header box
\setlength{\headsep}{0mm}     % Set the distance between the header box and the top of the text

\usepackage{gvv-book}
\usepackage{gvv}
\usepackage{cite}
\usepackage{amsmath,amssymb,amsfonts,amsthm}
\usepackage{algorithmic}
\usepackage{graphicx}
\graphicspath{{./figs/}}
\usepackage{textcomp}
\usepackage{xcolor}
\usepackage{txfonts}
\usepackage{listings}
\usepackage{enumitem}
\usepackage{mathtools}
\usepackage{gensymb}
\usepackage{comment}
\usepackage{float}
\usepackage[breaklinks=true]{hyperref}
\usepackage{tkz-euclide} 
\def\inputGnumericTable{}                                 
\usepackage[latin1]{inputenc}                                
\usepackage{color}                                            
\usepackage{array}                                            
\usepackage{longtable}                                       
\usepackage{calc}                                             
\usepackage{multirow}                                         
\usepackage{hhline}                                           
\usepackage{ifthen}                                           
\usepackage{lscape}
\usepackage{circuitikz}
\tikzstyle{block} = [rectangle, draw, fill=blue!20, 
text width=4em, text centered, rounded corners, minimum height=3em]
\tikzstyle{sum} = [draw, fill=blue!10, circle, minimum size=1cm, node distance=1.5cm]
\tikzstyle{input} = [coordinate]
\tikzstyle{output} = [coordinate]


\begin{document}
	
	\bibliographystyle{IEEEtran}
	\vspace{3cm}
	
	\title{2.7.28}
	\author{EE25BTECH11016 - Taraka Abhinav}
	\maketitle
	{\let\newpage\relax\maketitle}
	
	\setlength{\intextsep}{10pt} % Space between text and floats

\textbf{Question:}\\
\[
\text{Find the area of the triangle whose vertices are }
A(2,3),\; B(3,5),\; C(4,4).
\]

\solution \\
\begin{table}[H]    
  \centering
  \begin{center}
\begin{tabular}{|c|c|}
\hline
Point & Vector \\
\hline
$\vec{A}$ & $\begin{bmatrix} 2 \\ 3 \\ 0 \end{bmatrix}$ \\
\hline
$\vec{B}$ & $\begin{bmatrix} 3 \\ 5 \\ 0 \end{bmatrix}$ \\
\hline
$\vec{C}$ & $\begin{bmatrix} 4 \\ 4 \\ 0 \end{bmatrix}$ \\
\hline
\end{tabular}
\end{center}
  \caption{Vectors}
  \label{Answers}
\end{table}

\begin{align}
(\vec{A}-\vec{B})
&= \begin{bmatrix} 2 \\ 3 \\ 0 \end{bmatrix}
  - \begin{bmatrix} 3 \\ 5 \\ 0 \end{bmatrix}
= \begin{bmatrix} -1 \\ -2 \\ 0 \end{bmatrix},
\label{eq:AB} \\[6pt]
(\vec{A}-\vec{C})
&= \begin{bmatrix} 2 \\ 3 \\ 0 \end{bmatrix}
  - \begin{bmatrix} 4 \\ 4 \\ 0 \end{bmatrix}
= \begin{bmatrix} -2 \\ -1 \\ 0 \end{bmatrix}.
\label{eq:AC}
\end{align}

Using \eqref{eq:AB} and \eqref{eq:AC} The magnitude of the cross product is
\begin{align}
\big\lVert (\vec{A}-\vec{B}) \times (\vec{A}-\vec{C}) \big\rVert
&= \sqrt{0^2 + 0^2 + (-3)^2} = 3.
\label{eq:norm}
\end{align}

Therefore the area of triangle \(ABC\) is
\begin{align}
\text{ar}(\triangle ABC)
&= \tfrac{1}{2}\,\big\lVert (\vec{A}-\vec{B}) \times (\vec{A}-\vec{C}) \big\rVert
= \tfrac{1}{2}\times 3 = 1.5.
\label{eq:area}
\end{align}

      Therefore  \(\displaystyle \text{ar}(\triangle ABC)=1.5\).
\begin{figure}
    \centering
    \includegraphics[width=0.8\linewidth]{figs/Figure_1.png}
    \label{fig:placeholder}
\end{figure}
\end{document}